% ============================================================
% CAPITOLO 13 - FRASE E PERIODO
% ============================================================
\chapter{Frase e periodo}
\label{chap:struttura_frase}

Gli accordi di un brano formano un discorso organizzato, che le cadenze dividono in unità strutturali: le frasi e i periodi.

\section{La fraseologia}
\label{sec:fraseologia}

Le cadenze sono i punti di arrivo di un percorso governato dalle \textbf{funzioni armoniche} (Capitolo \ref{chap:funzioni_armoniche}). La struttura più comune di una frase si riassume nel modello T-S-D-T.

\begin{regola}[Logica funzionale: T-S-D-T]
	\label{reg:tsdt}
	Il modello \textbf{T-S-D-T} descrive il ciclo di tensione e risoluzione alla base della sintassi tonale:
	\begin{itemize}
		\item \textbf{T (tonica):} area di stabilità e riposo (I grado).
		\item \textbf{S (sottodominante):} area di allontanamento dalla tonica (IV e II grado).
		\item \textbf{D (dominante):} area di massima tensione, che prepara il ritorno alla tonica (V grado).
		\item \textbf{T (tonica):} ritorno al riposo, chiusura del ciclo.
	\end{itemize}
	Su questo ciclo è costruita l'armonia di moltissimi brani.
\end{regola}

\begin{esempio}[Una frase basata sul modello T-S-D-T]
	\label{ex:tsdt}
	Una realizzazione semplice di I - IV - V - I.
	\musica{es_cap06_07}
\end{esempio}

\section{La frase musicale}
\label{sec:frase_musicale}

\begin{definizione}[Frase musicale]
    È l'unità minima di senso compiuto nel discorso tonale. È composta tipicamente di 4 misure (con varianti) e termina con una cadenza che ne segna il confine.
\end{definizione}

Come una frase verbale ha un soggetto, uno sviluppo e un punto finale, la frase musicale ha un inizio, un percorso armonico e una conclusione cadenzale. Per questo, in un basso dato, ogni frase si deve chiudere con una cadenza riconoscibile.

\section{Il periodo}
\label{sec:periodo}

\begin{definizione}[Il periodo]
    È una struttura di norma di 8 misure, formata da due frasi in relazione fra loro:
    \begin{itemize}
        \item \textbf{antecedente} (mis. 1-4): termina con una cadenza non conclusiva, sospesa o imperfetta (Definizioni \ref{def:cadenza_sospesa} e \ref{def:cadenza_imperfetta});
        \item \textbf{conseguente} (mis. 5-8): termina con una cadenza conclusiva, la perfetta (Definizione \ref{def:cadenza_perfetta}).
    \end{itemize}
\end{definizione}

L'antecedente e il conseguente stanno come domanda e risposta: la prima frase crea un'aspettativa e la seconda la soddisfa.

\begin{center}
\begin{tabular}{lcl}
\toprule
\textbf{Frase} & \textbf{Misure} & \textbf{Cadenza} \\
\midrule
Antecedente & 1--4 & sospesa su V, oppure imperfetta \\
Conseguente & 5--8 & perfetta V-I \\
\bottomrule
\end{tabular}
\end{center}

\section{Pianificare le cadenze nel basso dato}
\label{sec:applicazione_frase}

\begin{regola}[Pianificazione cadenzale]
    Prima di armonizzare un basso si individuano:
    \begin{itemize}
        \item \textbf{metà (mis. 4 di 8):} una cadenza sospesa (V) o imperfetta (I non fondamentale, oppure soprano non sulla tonica);
        \item \textbf{finale (mis. 8):} una cadenza perfetta (V-I in stato fondamentale).
    \end{itemize}
\end{regola}

\begin{regola}[Scelta della cadenza di metà]
    Si preferisce la cadenza sospesa su V se il basso si ferma sulla dominante, e la cadenza imperfetta se il basso approda al I ma senza la forza conclusiva necessaria per il finale.
\end{regola}
